\begin{table}[!htbp]
\centering\small
\setlength{\tabcolsep}{4pt}
\renewcommand{\arraystretch}{1.12}
\begin{tabularx}{\linewidth}{@{}>{\raggedright\arraybackslash}p{.16\linewidth}>{\raggedright\arraybackslash}X>{\raggedright\arraybackslash}X@{}}
\toprule
Grandeur & Résultat dans la réalisation spécifiée & Référence et comparaison\\
\midrule
Planck, avant le retour
& \(\hbar_{\rm pre}=1.0545718176461544696\ldots\)
  \(\times10^{-34}\,\mathrm{J\,s}\).
& SI : \(\hbar_{\rm SI}=h_{\rm SI}/(2\pi)\), exact.
  Sa coordonnée est \(1.0545718176461563912\ldots\)
  \(\times10^{-34}\,\mathrm{J\,s}\).
  Résidu relatif \(-1.822221\times10^{-15}\).\\[4pt]
Planck, après le retour
& \(\hbar_{\rm ret}=1.0545718169150390611\ldots\)
  \(\times10^{-34}\,\mathrm{J\,s}\).
& Même référence SI.
  Résidu relatif \(-6.932836\times10^{-10}\).
  Pour \(h_a=2\pi\hbar_a\), les résidus sont identiques.\\[4pt]
Barbero--Immirzi
& \(\gamma_{\HMT}=0.2740076726944592747\ldots\).
& Ghosh--Mitra : environ \(0.2740668582328300\).
  Différence relative \(-2.159529202\times10^{-4}\)
  entre constructions théoriques.\\[4pt]
Mélange CKM
& \(\delta_{\rm CP}=65.676173123554^\circ\),
  \(J=3.1189723326632974\times10^{-5}\).
  Les trois angles de \eqref{eq:ii-ckm-decimales} sont
  \(\theta_{12}=13.003313589509^\circ\),
  \(\theta_{23}=2.398056157332^\circ\) et
  \(\theta_{13}=0.213977382244^\circ\).
& PDG 2025 : les cinq comparaisons de
  \(s_{12},s_{23},s_{13},\delta,J\) satisfont \(|z_x|<0.029\)
  avec les incertitudes marginales du
  tableau~\ref{tab:ii-ckm-comparacion}.\\[4pt]
Boltzmann
& \(B_*=1.380649\times10^{-23}\,u_E/u_\Theta\),
  avec les bases et la règle thermométrique de
  \eqref{eq:ii-b-evaluacion-comparada}.
& SI : \(1.380649\times10^{-23}\,\mathrm{J/K}\), exact.
  Coïncidence du nombre sous le transport dimensionnel déclaré.\\[4pt]
Longueur \(L\)
& \(L=1.6162549826251263301\ldots\)
  \(\times10^{-35}\,\mathrm m\), pour \(L_*=1\,\mathrm m\).
& Composition longitudinale et radiale :
  \(G_{\rm ret}=c_{\rm int}^3L^2/\hbar_{\rm ret}\),
  avec les mêmes bases de vitesse et d’action.\\[4pt]
Gravitation
& \(G_{\rm ret}=6.6742996594140227\ldots\)
  \(\times10^{-11}\,\mathrm{m^3\,kg^{-1}\,s^{-2}}\).
& CODATA 2022 : \(6.67430(15)\times10^{-11}\) dans les mêmes unités.
  Différence d’environ \(-0.00227057\) incertitudes-types.\\
\bottomrule
\end{tabularx}
\caption{Synthèse de la comparaison quantitative. Les évaluations, les bases
et les références sont celles des tableaux~\ref{tab:ii-planck-comparacion}--\ref{tab:ii-kb-g-cartas}.
Les résidus relatifs des constantes d’action, la comparaison
théorique de Barbero--Immirzi et les résidus normalisés de CKM et de la gravitation
conservent leurs sens respectifs.}
\label{tab:ii-conclusiones-contraste}
\par\medskip
\centering\small
\setlength{\tabcolsep}{4pt}
\renewcommand{\arraystretch}{1.18}
\begin{tabular}{@{}lrrr@{}}
\toprule
Observable & Évaluation HMT & PDG 2025 & \(z_x\)\\
\midrule
\(s_{12}\) & \(0.225007404757\) & \(0.22501\pm0.00068\) & \(-0.003817\)\\
\(s_{23}\) & \(0.041841757010\) & \(0.04183^{+0.00079}_{-0.00069}\) & \(0.014882\)\\
\(s_{13}\) & \(0.003734601164\) & \(0.003732^{+0.000090}_{-0.000085}\) & \(0.028902\)\\
\(\delta\) (rad) & \(1.146265461116\) & \(1.147\pm0.026\) & \(-0.028251\)\\
\(10^5J\) & \(3.118972333\) & \(3.12^{+0.13}_{-0.12}\) & \(-0.008564\)\\
\bottomrule
\end{tabular}
\caption{Valeurs CKM et comparaison marginale réunies dans les conclusions.
Les trois angles du tableau~\ref{tab:ii-conclusiones-contraste} sont comparés
au moyen de \(s_{ij}=\sin\theta_{ij}\), avec la phase en radians et
l’invariant \(J\). On conserve les chiffres, incertitudes et résidus du
tableau~\ref{tab:ii-ckm-comparacion}, correspondant à l’ajustement à trois
générations PDG 2025. La dernière ligne exprime conjointement la valeur et
l’incertitude en unités de \(10^{-5}\). Les résidus sont des comparaisons
marginales d’observables corrélées.}
\label{tab:ii-conclusiones-ckm}
\end{table}
